% !Mode:: "TeX:UTF-8"

\chapter{绪论}

\section{研究背景与意义}

\subsection{研究背景}

这里填写研究背景。建议交代现实问题、研究对象与学术背景，并结合正式文献说明选题来源\cite{sample_article,sample_book}。

\subsection{研究意义}

这里填写理论意义与实践意义。写作时建议围绕“问题在哪里、本文解决什么、贡献体现在哪里”展开，避免空泛套话。

\section{国内外研究现状}

这里综述相关研究。建议按主题或方法分组，而不是简单堆砌作者与年份；同时确保文中表述与参考文献条目一致。

\section{本文主要工作与论文结构}

这里概括本文的主要工作，并说明各章安排。若学校或学院要求“绪论”包含研究思路、技术路线或创新点，也可在本节后补充小节。

\chapter{相关技术与理论基础}

\section{核心概念}

这里说明本文涉及的核心概念、定义与理论基础。首次出现的重要缩写建议采用“中文名称（English Name, 缩写）”的形式。

\section{关键技术或分析框架}

这里介绍后文需要用到的方法、模型、指标或理论框架，并说明它们与本文研究问题的关系。

\chapter{系统设计或研究方法}

\section{总体设计}

这里描述系统架构、研究框架或分析流程。若使用图示，请确保图中术语、层级与正文叙述严格一致。

\section{核心机制}

这里展开方法设计、算法流程或分析路径。对于公式、变量和参数，应在首次出现时说明其含义与取值范围。

\begin{equation}
  y = f(x)
\end{equation}

\section{示例表格}

表格建议优先使用三线表，表题放在表上方，表内指标和单位应完整。

\begin{table}[H]
  \centering
  \wuhao
  \caption{示例三线表}
  \label{tab:sample}
  \begin{tabular}{lcc}
    \toprule
    指标 & 方案A & 方案B \\
    \midrule
    指标一 & 0.80 & 0.85 \\
    指标二 & 0.72 & 0.79 \\
    \bottomrule
  \end{tabular}
\end{table}

\chapter{实验结果或案例分析}

\section{实验设置或案例背景}

这里写数据来源、样本范围、实验环境、评价指标或案例背景。若涉及多组数据，建议明确年份、口径和可比性边界。

\section{结果分析}

这里结合图表说明结果。正文应解释“结果说明了什么”和“为什么会出现这样的结果”，而不是只重复表格中的数值。

\section{讨论}

这里补充局限性、误差来源、边界条件或与已有研究的差异。

\chapter{总结与展望}

\section{工作总结}

这里概括全文完成的主要工作和得到的核心结论，注意语言收束，不要重复大段背景。

\section{不足与展望}

这里说明本文的不足以及未来可以继续推进的方向。
